\documentclass[12pt,a4paper]{article}
\usepackage[utf8]{inputenc}
\usepackage[T1]{fontenc}
\usepackage{enumitem}
\usepackage[margin=2.5cm]{geometry}
\usepackage{parskip}
\usepackage{amsmath}
\usepackage{amssymb}

\title{midterm --- Final Exam}
\date{}

\begin{document}

\maketitle



\noindent
\textbf{Total points:} 4
\hfill
\textbf{Number of questions:} 4

\vspace{1em}
\hrule
\vspace{1.5em}


\noindent
\textbf{Question 1 (1.00 pts).}~
What is a key difference between a classical bit and a qubit?


\begin{enumerate}[label=\textbf{\Alph*.}]

  \item A classical bit can be in a blend of states at once, while a qubit cannot.

  \item A qubit is like a light switch, always in exactly one of two positions.

  \item A qubit is more like a coin spinning in the air, with an undecided outcome until measured.

  \item A classical bit is more like a coin spinning in the air, with an undecided outcome until measured.

\end{enumerate}



\vspace{0.5em}


\noindent
\textbf{Question 2 (1.00 pts).}~
Quantum gates can destroy or create probability.


\begin{enumerate}[label=\textbf{\Alph*.}]
  \item True
  \item False
\end{enumerate}



\vspace{0.5em}


\noindent
\textbf{Question 3 (1.00 pts).}~
What is the key difference between the correlation of entangled qubits and the correlation of matching socks?


\vspace{0.4cm}
\noindent\makebox[\linewidth]{\hrulefill}
\vspace{0.6cm}



\vspace{0.5em}


\noindent
\textbf{Question 4 (1.00 pts).}~
Explain the concept of superposition in quantum computing and how it differs from classical computing.


\vspace{0.4cm}
\noindent\makebox[\linewidth]{\hrulefill}
\vspace{0.4cm}
\noindent\makebox[\linewidth]{\hrulefill}
\vspace{0.4cm}
\noindent\makebox[\linewidth]{\hrulefill}
\vspace{0.6cm}





\end{document}
